% Expressions for Paradiso's section locations $\bm{r}_*$ in the notation of draft.tex
\newcommand{\WarpTwistCentroid}{\bm{r}_{wtc}}
\newcommand{\ShearForceLocation}{\bm{\rho}_{F}}
\subsection*{Setup and Assumptions}

Work in centroidal coordinates so that the transverse origin is at the centroid:
\begin{equation}
  \CentroidLocation = \bar{\bm{\SectionLocation{}}} = \mathbf{0},
  \qquad
  \SectionCoordinate = \bm{r},
  \qquad
  \IntRGoRG = \mathbf{J}_{\!G} = \int_\Section \bm{r}\otimes\bm{r}\,\mathrm{d}A.
\end{equation}
If a non-centroidal frame is used, replace every occurrence of $\bm{r}$ by $\bm{r} - \CentroidLocation$.

All symbols such as $\AxialStress$, $\ShearStress$, $\ShearForce$, $\SectionCoordinate$, $\WarpTwistIesan$, $\WarpShearIesan$, $\BishearTensorPB$, $\IntRGoRG$, $\FrameBasis$, $\CenterShearLocation$, $\CenterTwistLocation$ etc.\ are those already defined in \texttt{draft.tex}.

\subsection*{1.\ $\bm{r}_F$ (center of normal stress)}

In Paradiso, for the axial stress field $\sigma_{\mathrm{ax}}$ and axial resultant $N$,
\begin{equation}
  \int_\Section \sigma_{\mathrm{ax}}\,\mathrm{d}A = N,
  \qquad
  \int_\Section \sigma_{\mathrm{ax}}\,\bm{r}\,\mathrm{d}A = N\,\bm{r}_F.
\end{equation}
In the notation of \texttt{draft.tex}, with $\AxialStress$ the axial stress and $\AxialRxn$ the axial resultant,
\begin{equation}
  \AxialRxn := \int_\Section \AxialStress\,\mathrm{d}A,
  \qquad
  \AxialRxn\,\bm{\SectionLocation}_F := \int_\Section \AxialStress\,\SectionCoordinate\,\mathrm{d}A.
\end{equation}
Thus
\begin{equation}
  \boxed{\bm{r}_F \;\longleftrightarrow\; \bm{\SectionLocation}_F
  \text{ defined by }
  \AxialRxn\,\bm{\SectionLocation}_F = \int_\Section \AxialStress\,\SectionCoordinate\,\mathrm{d}A.}
\end{equation}

\subsection*{2.\ $\bm{r}_P$ (prescribed load point)}

Paradiso defines $\bm{r}_P$ as the position in the cross-section where the shear force is applied. This is prescribed data, not determined by section integrals.

In \texttt{draft.tex}, the corresponding object is the \emph{trace center} $\bm{r}_{\mathcal T}$, i.e.\ the section location at which the trace line and FE node are placed:
\begin{equation}
  \boxed{\bm{r}_P \;\longleftrightarrow\; \bm{r}_{\mathcal T}.}
\end{equation}

\subsection*{3.\ $\bm{r}_p$ (statical equivalence point for the shear field)}

Paradiso defines $\bm{r}_p$ by the statical equivalence conditions for the shear stress field $\bm{\tau}_{\mathrm{sh}}$ associated with a shear force $\ShearForce$:
\begin{equation}
  \int_\Section \bm{\tau}_{\mathrm{sh}}\,\mathrm{d}A = \ShearForce,
\end{equation}
\begin{equation}
  \int_\Section \bm{r}\times\bm{\tau}_{\mathrm{sh}}\cdot\FrameBasis\,\mathrm{d}A
  = (\bm{r}_p\times\ShearForce)\cdot\FrameBasis.
\end{equation}

In \texttt{draft.tex} notation, with $\ShearStress$ the shear stress and $\ShearForce$ the shear resultant,
\begin{equation}
  \ShearForce := \int_\Section \ShearStress\,\mathrm{d}A,
\end{equation}
\begin{equation}
  (\bm{\SectionLocation}_p\times\ShearForce)\cdot\FrameBasis
  := \int_\Section \bigl(\SectionCoordinate\times\ShearStress\bigr)\cdot\FrameBasis\,\mathrm{d}A.
\end{equation}
Thus
\begin{equation}
  \boxed{\bm{r}_p \;\longleftrightarrow\; \bm{\SectionLocation}_p
  \text{ defined by }
  (\bm{\SectionLocation}_p\times\ShearForce)\cdot\FrameBasis
  = \int_\Section (\SectionCoordinate\times\ShearStress)\cdot\FrameBasis\,\mathrm{d}A.}
\end{equation}

\subsection*{4.\ $\bm{r}_T$ (center of twist)}

Paradiso defines the center of twist $\bm{r}_T$ by
\begin{equation}
  \bm{r}_T := \FrameBasis\times\mathbf{J}_G^{-1}\,\mathbf{s}_{\WarpTwistCentroid},
  \qquad
  \mathbf{s}_{\WarpTwistCentroid} := \int_\Section \WarpTwistCentroid\,\bm{r}\,\mathrm{d}A,
\end{equation}
where $\WarpTwistCentroid$ is the torsion warping scalar and $\mathbf{J}_G$ is the centroidal inertia tensor.

In \texttt{draft.tex}, with $\WarpTwistIesan$ the torsion warping in the Ieşan representation and $\IntRGoRG = \mathbf{J}_G$,
\begin{equation}
  \mathbf{s}_{\WarpTwistIesan} := \int_\Section \WarpTwistIesan\,\SectionCoordinate\,\mathrm{d}A,
\end{equation}
and the corresponding twist center is
\begin{equation}
  \boxed{
  \CenterTwistLocation
  = \FrameBasis\times\IntRGoRG^{-1}\,\mathbf{s}_{\WarpTwistIesan}
  = \FrameBasis\times\IntRGoRG^{-1}
    \int_\Section \WarpTwistIesan\,\SectionCoordinate\,\mathrm{d}A.}
\end{equation}
Thus
\begin{equation}
  \bm{r}_T \;\longleftrightarrow\; \CenterTwistLocation.
\end{equation}

\subsection*{5.\ $\bm{r}_\psi$ (shear/geometric center associated with shear warping)}

Paradiso defines a shear-related center $\bm{r}_\psi$ by
\begin{equation}
  \bm{r}_\psi
  = \FrameBasis\times\mathbf{J}_G^{-1}\,\bm{\gamma}_\psi,
  \qquad
  \bm{\gamma}_\psi
  = \int_\Section \mathbf{H}_\psi\bigl(\FrameBasis\times\bm{r}\bigr)\,\mathrm{d}A,
\end{equation}
where $\mathbf{H}_\psi := \bm{\psi}\otimes\nabla + \mathbf{A}^p$ is the shear warping operator built from the harmonic shear warping $\bm{\psi}$ and the isotropic particular solution $\mathbf{A}^p$.

In \texttt{draft.tex}, with
\begin{equation}
  \WarpShearIesan = \psi,
  \qquad
  \BishearTensorPB = \mathbf{H}_\psi,
\end{equation}
and again $\IntRGoRG = \mathbf{J}_G$, the corresponding center is
\begin{equation}
  \bm{\gamma}_\psi
  = \int_\Section \BishearTensorPB\bigl(\FrameBasis\times\SectionCoordinate\bigr)\,\mathrm{d}A,
\end{equation}
\begin{equation}
  \boxed{
  \CenterShearLocation
  = \FrameBasis\times\IntRGoRG^{-1}\,\bm{\gamma}_\psi
  = \FrameBasis\times\IntRGoRG^{-1}
    \int_\Section \BishearTensorPB\bigl(\FrameBasis\times\SectionCoordinate\bigr)\,\mathrm{d}A.}
\end{equation}
Thus
\begin{equation}
  \bm{r}_\psi \;\longleftrightarrow\; \CenterShearLocation.
\end{equation}

\subsection*{6.\ $\bm{r}_{\psi_p}$ (shear-warping shift relative to $\bm{r}_\psi$)}

Paradiso defines
\begin{equation}
  \bm{r}_{\psi_p} := \bm{r}_p - \bm{r}_\psi.
\end{equation}
Using the identifications above,
\begin{equation}
  \bm{r}_p \;\longleftrightarrow\; \bm{\SectionLocation}_p,
  \qquad
  \bm{r}_\psi \;\longleftrightarrow\; \CenterShearLocation,
\end{equation}
the corresponding vector in the notation of \texttt{draft.tex} is
\begin{equation}
  \boxed{
  \bm{\SectionLocation}_{\WarpShearIesan,p}
  := \bm{\SectionLocation}_p - \CenterShearLocation.}
\end{equation}
Thus
\begin{equation}
  \bm{r}_{\psi_p} \;\longleftrightarrow\; \bm{\SectionLocation}_p - \CenterShearLocation.
\end{equation}

\subsection*{7.\ Section locations not explicitly expressed}

The following objects from Paradiso are not expressed here purely in terms of the symbols of \texttt{draft.tex}, in order to avoid incomplete or speculative formulas:
\begin{itemize}
  \item The explicit integral expression for $\bm{\gamma}_\psi$ in terms of the section tensors
        $\mathbf{H}_{\psi\psi}$, $\mathbf{h}_{\varphi\psi}$, $J_0$ as they appear in Sec.~A.2.
  \item A closed-form expression for $\bm{r}_p$ solely in terms of those section tensors, obtained by
        substituting the explicit stress field into the statical equivalence conditions and solving.
\end{itemize}
